% =====================================================================
\section{Phương pháp đề xuất}
\label{sec:method}
% =====================================================================

\subsection{Kiến trúc học đa nhiệm}
\label{sec:arch}

Backbone là ResNet-18 \textbf{khởi tạo ngẫu nhiên}, không dùng trọng số tiền huấn luyện. Việc
không dùng tiền huấn luyện là lựa chọn có chủ đích: nếu encoder đã được huấn luyện trước trên
ImageNet, hiệu ứng của nhiệm vụ phụ sẽ bị trộn lẫn với chất lượng đặc trưng sẵn có và không
còn đo được đóng góp của tín hiệu tự giám sát~\cite{he2016resnet,pytorch_resnet}.

Điều chỉnh cho ảnh $32 \times 32$:
\begin{itemize}[leftmargin=1.2cm,itemsep=2pt]
  \item thay convolution đầu bằng \code{Conv2d(3, 64, kernel\_size=3, stride=1, padding=1, bias=False)};
  \item bỏ lớp MaxPool đầu mạng, vì sau hai lần giảm độ phân giải $32 \to 8$ thì thông tin
        không gian còn quá ít;
  \item giữ bốn nhóm residual block, mỗi nhóm hai BasicBlock (cấu hình ResNet-18 chuẩn).
\end{itemize}

% Sinh tự động — scripts/make_report_data.py
\begin{table}[htbp]
  \centering
  \small
  \begin{tabular}{@{}l l@{}}
    \toprule
    Thành phần & Kích thước đầu ra $C \times H \times W$ \\
    \midrule
    Ảnh đầu vào & $3 \times 32 \times 32$ \\
    Conv $3\times3$ + BatchNorm + ReLU & $64 \times 32 \times 32$ \\
    Nhóm residual 1: 2 BasicBlock & $64 \times 32 \times 32$ \\
    Nhóm residual 2: 2 BasicBlock & $128 \times 16 \times 16$ \\
    Nhóm residual 3: 2 BasicBlock & $256 \times 8 \times 8$ \\
    Nhóm residual 4: 2 BasicBlock & $512 \times 4 \times 4$ \\
    Global Average Pooling + Flatten & vector 512 chiều \\
    Superclass Head (tuyến tính, có bias) & $512 \to 2$ \\
    Rotation Head (tuyến tính, có bias) & $512 \to 4$ \\
    \bottomrule
  \end{tabular}
  \caption{Chuỗi kích thước tensor của backbone ResNet-18 điều chỉnh cho ảnh $32 \times 32$. Các giá trị này được kiểm tra tự động bằng \texttt{scripts/self\_test.py} trên môi trường thực nghiệm.}
  \label{tab:arch}
\end{table}

\paragraph{Ký hiệu.}
$\enc$ là encoder dùng chung; $\gsup$ và $\grot$ là hai head tương ứng; $\norm$ là bước chuyển
tensor và chuẩn hóa; $R_k$ là phép xoay $90k$ độ. Hai đầu ra được viết:

\begin{align}
  a_{\mathrm{sup}} &= \gsup\bigl(\enc(t_{\mathrm{sup}}(x))\bigr) \in \mathbb{R}^2, \label{eq:asup}\\
  a_{\mathrm{rot}} &= \grot\bigl(\enc(\norm(R_k(x)))\bigr) \in \mathbb{R}^4, \label{eq:arot}
\end{align}

trong đó $t_{\mathrm{sup}}(\cdot)$ là chuỗi augmentation của nhánh phân loại. Việc viết tách
$t_{\mathrm{sup}}$ và $\norm \circ R_k$ nhấn mạnh rằng \textbf{hai nhánh nhận các view khác
nhau}; ta không giả định tồn tại một vector đặc trưng duy nhất dùng cho cả ảnh chuẩn và ảnh đã
xoay.

\paragraph{Cách ghép minibatch.}
Khi huấn luyện M8, hai minibatch được ghép theo chiều batch trước khi qua encoder:
$[\,B_L;\, B_R\,] \in \mathbb{R}^{(b_L + b_R) \times 3 \times 32 \times 32}$, sau đó tách đặc
trưng trở lại đúng head. \textbf{Các lớp BatchNorm dùng chung}, do đó mọi thay đổi về thành phần
của batch đều làm thay đổi thống kê chạy của toàn mạng~\cite{ioffe2015batchnorm}. Đây là cơ sở
cho một giới hạn quan trọng được nêu ở \secref{sec:exp-compare}: không thể diễn giải so sánh
\cfg{E2-L}/\cfg{E1} như một phép tách hoàn toàn riêng tác dụng của nhãn góc xoay khỏi mọi thay
đổi về mặt tính toán.

\paragraph{Suy luận.}
Khi suy luận, chỉ chạy encoder và Superclass Head; chọn lớp bằng $\arg\max$ của hai logits.
Rotation Head không cần thiết và không được tính vào chi phí suy luận. Hệ quả kiến trúc: mô
hình suy luận của \cfg{E2} có \emph{đúng cùng} kiến trúc với \cfg{E1}, nên không thể kỳ vọng
giảm độ phức tạp kiến trúc chỉ nhờ bỏ head phụ.

\subsection{Hàm mất mát và cách sử dụng dữ liệu}
\label{sec:loss}

Với $B_L$ là minibatch có nhãn, $B_R$ là minibatch cho nhiệm vụ xoay, và $b_L$, $b_R$ là số mẫu
tương ứng, hai mất mát được lấy \textbf{trung bình riêng}:

\begin{align}
  \Lsup &= \frac{1}{b_L}\sum_{i=1}^{b_L} \mathrm{CE}\bigl(a_{\mathrm{sup},i},\, y_i\bigr),
  \label{eq:lsup}\\
  \Lrot &= \frac{1}{b_R}\sum_{j=1}^{b_R} \mathrm{CE}\bigl(a_{\mathrm{rot},j},\, k_j\bigr),
  \label{eq:lrot}\\
  L &= \Lsup + \lamrot \Lrot. \label{eq:total}
\end{align}

$\mathrm{CE}$ là cross-entropy cho nhãn lớp nguyên. Trong PyTorch, logits được truyền trực tiếp
vào \code{CrossEntropyLoss}, \textbf{không} áp dụng Softmax trước loss; việc áp Softmax hai lần
là một lỗi phổ biến làm hỏng gradient~\cite{pytorch_ce}. Trong mã nguồn, logits được truyền
nguyên vẹn vào lớp này.

\paragraph{Phân chia nguồn dữ liệu cho từng mất mát.}
Chỉ ảnh thuộc $\DL$ đóng góp vào $\Lsup$. Ngược lại, $B_R$ lấy từ $\DL$ hoặc từ $\Dtrain$ tùy
cấu hình (\tabref{tab:configs}). Chốt $\lamrot = 0{,}5$ cho thí nghiệm chính; các giá trị khác
chỉ thuộc khảo sát độ nhạy đã khai báo trước.

\begin{keybox}
\textbf{Vì sao lấy trung bình riêng từng mất mát?} Nếu gộp hai tập mẫu rồi lấy trung bình chung,
trọng số hiệu dụng của nhánh xoay sẽ phụ thuộc vào tỷ lệ $b_R/b_L$, và như vậy $\lamrot$ mất ý
nghĩa so sánh giữa các mức nhãn. Lấy trung bình riêng giữ cho $\lamrot$ là một hệ số có thể
diễn giải được và giữ phép so sánh độ nhạy $\lamrot$ nhất quán.
\end{keybox}

\subsection{Mở rộng Contrastive Learning và SupCon}
\label{sec:contrastive}

Nội dung này là \textbf{hướng mở rộng} của đề tài: nó không phải điều kiện hoàn thành phần
chính và được chạy \emph{sau} khi khối đối chứng \cfg{E1}/\cfg{E2-L}/\cfg{E2} đã hoàn tất.
Bốn cấu hình \cfg{E3}--\cfg{E6} được triển khai và báo cáo tại \secref{sec:results-ext}.

\paragraph{Contrastive Learning tự giám sát (NT-Xent).}
Tạo hai view của cùng một ảnh và đưa qua encoder cùng Projection Head; hai view này tạo thành
positive pair. Có thể dùng ảnh từ $\Dtrain$ vì \textbf{không} cần nhãn siêu lớp. Theo tham chiếu
SimCLR~\cite{chen2020simclr}, mọi thành phần đều được khóa trước: loss NT-Xent, chuẩn hóa vector
$L_2$, nhiệt độ $\tau_{\mathrm{NT}}$, phép tăng cường ảnh, kích thước batch và chiều Projection
Head. Với $2N$ view trong một minibatch và $j(i)$ là view còn lại của cùng ảnh:

\begin{equation}
  L_{\mathrm{NT\text{-}Xent}}
  = \frac{1}{2N}\sum_{i=1}^{2N} -\log
    \frac{\exp\!\bigl(s_{i,j(i)}/\tau_{\mathrm{NT}}\bigr)}
         {\sum_{k \neq i} \exp\!\bigl(s_{i,k}/\tau_{\mathrm{NT}}\bigr)},
  \qquad
  s_{i,k} = \frac{z_i^{\top} z_k}{\lVert z_i\rVert_2 \lVert z_k\rVert_2}.
  \label{eq:ntxent}
\end{equation}

\paragraph{Supervised Contrastive Learning (SupCon).}
Xây dựng positive pair từ \textbf{các view có cùng nhãn siêu lớp đã được cấp cho huấn luyện}.
Trong đề tài này, SupCon chỉ sử dụng $\DL$; \textbf{không} đọc nhãn thật của $\DU$. SupCon là
học \emph{có} giám sát, không được gọi là một nhiệm vụ tự giám sát~\cite{khosla2020supcon}. Với
$A(i) = \{a \neq i : \tilde{y}_a = \tilde{y}_i\}$ là tập positive của view $i$:

\begin{equation}
  L_{\mathrm{SupCon}}
  = \sum_{i=1}^{2N} \frac{-1}{|A(i)|}
    \sum_{a \in A(i)} \log
    \frac{\exp\!\bigl(s_{i,a}/\tau_{\mathrm{Sup}}\bigr)}
         {\sum_{k \neq i} \exp\!\bigl(s_{i,k}/\tau_{\mathrm{Sup}}\bigr)}.
  \label{eq:supcon}
\end{equation}

\paragraph{Kiến trúc Projection Head.}
Hai lớp tuyến tính $512 \to 512 \to \dim_{\mathrm{proj}}$ với BatchNorm và ReLU ở lớp ẩn
(theo SupCon/SimCLR). Head này \textbf{chỉ} tồn tại khi huấn luyện; khi suy luận, mô hình chỉ
gồm encoder và Superclass Head, nên \emph{số tham số suy luận không thay đổi} so với
\cfg{E1}/\cfg{E2} (xem \tabref{tab:inference-cost} và \tabref{tab:extension-cost}).

\subsection{Đối chứng bắt buộc}
\label{sec:baselines}

Ba cấu hình chính được thiết kế để tách bạch hai yếu tố thường bị trộn lẫn: \emph{việc thêm
nhiệm vụ phụ} và \emph{việc mở rộng nguồn ảnh cho nhiệm vụ phụ}.

% Sinh tự động — scripts/make_report_data.py
\begin{table}[htbp]
  \centering
  \small
  \begin{tabular}{@{}l l l l@{}}
    \toprule
    ID & Cấu hình & Ảnh cho phân loại & Ảnh cho nhiệm vụ xoay \\
    \midrule
    \cfg{E1}   & Baseline chỉ học siêu lớp & $\DL$ & Không có \\
    \cfg{E2-L} & Thêm nhánh xoay, cùng tập nguồn & $\DL$ & Chỉ $\DL$ \\
    \cfg{E2}   & Thêm nhánh xoay, mở rộng nguồn ảnh & $\DL$ & $\Dtrain = \DL \cup \DU$ \\
    \bottomrule
  \end{tabular}
  \caption{Ba cấu hình chính. \cfg{E1} và \cfg{E2} là cặp đối chứng chính; \cfg{E2-L} tách riêng tác dụng của việc mở rộng nguồn ảnh cho nhánh phụ.}
  \label{tab:configs}
\end{table}

% Sinh tự động — scripts/make_report_data.py
\begin{table}[htbp]
  \centering
  \footnotesize
  \setlength{\tabcolsep}{4pt}
  \begin{tabular}{@{}l l l l@{}}
    \toprule
    ID & Hàm mục tiêu & Nguồn ảnh cho xoay & Nguồn ảnh cho loss bổ sung \\
    \midrule
    \cfg{E3} & $\Lsup + \lamc \Lntxent$ & Không có nhánh xoay & $\Dtrain$; không dùng nhãn siêu lớp \\
    \cfg{E4} & $\Lsup + \lamc \Lsupcon$ & Không có nhánh xoay & Chỉ $\DL$ có nhãn hợp lệ \\
    \cfg{E5} & $\Lsup + \lamrot \Lrot + \lamc \Lntxent$ & $\Dtrain$ & NT-Xent: $\Dtrain$; rotation: $\Dtrain$ \\
    \cfg{E6} & $\Lsup + \lamrot \Lrot + \lamc \Lsupcon$ & $\Dtrain$ & SupCon: $\DL$; rotation: $\Dtrain$ \\
    \bottomrule
  \end{tabular}
  \caption{Bốn cấu hình mở rộng contrastive. $\lamc$ dùng chung cho mọi cấu hình và được khóa bằng validation trước khi đánh giá test (\secref{sec:results-ext-lambda}). SupCon chỉ dùng $\DL$; NT-Xent dùng $\Dtrain$ vì không cần nhãn siêu lớp.}
  \label{tab:extensions}
\end{table}

Ba so sánh ghép cặp được dùng để trả lời các câu hỏi đã nêu:

\begin{itemize}[leftmargin=1.2cm,itemsep=3pt]
  \item \textbf{\cfg{E2-L} so với \cfg{E1}:} tác động của việc bổ sung nhánh rotation khi
        \emph{không} mở rộng tập ảnh nguồn. Đây là phép so sánh cô lập nhất về mặt dữ liệu,
        nhưng vẫn còn khác biệt về tính toán.
  \item \textbf{\cfg{E2} so với \cfg{E2-L}:} ảnh hưởng của việc mở rộng nguồn ảnh cho nhánh phụ
        (chỉ $\DL$ so với toàn bộ $\Dtrain$). Hai cấu hình này có cùng hàm mục tiêu và cùng số
        bước cập nhật, chỉ khác nguồn ảnh của $B_R$.
  \item \textbf{\cfg{E2} so với \cfg{E1}:} hiệu quả tổng thể của M8 trong cùng ngân sách nhãn;
        đây là so sánh trả lời trực tiếp câu hỏi C1.
\end{itemize}

Các so sánh được ghép cặp theo lần lặp và tỷ lệ nhãn, tức $\Delta_s = \mathrm{Macro\text{-}F1}_s(\text{M8}) -
\mathrm{Macro\text{-}F1}_s(\cfg{E1})$ tính riêng cho từng $s$ rồi mới tổng hợp.

\begin{warnbox}
\textbf{Giới hạn về tính công bằng của phép so sánh.} \cfg{E2-L}/\cfg{E2} vẫn xử lý thêm view và
làm thay đổi thống kê BatchNorm so với \cfg{E1}. Do đó \textbf{không} diễn giải chúng là phép
tách hoàn toàn riêng tác dụng của \enquote{nhãn góc xoay} khỏi mọi thay đổi về tính toán.
Tương tự, \textbf{không} khẳng định so sánh công bằng về FLOPs hoặc thời gian chỉ vì ba cấu hình
dùng cùng số epoch: số lượt ảnh được xử lý ở mỗi bước cập nhật là khác nhau.
\end{warnbox}
